% Programa Sanepar Rural: sequência de atividades por responsável (raias)
\begin{tikzpicture}[passo/.style={etapa=2.2cm, minimum height=1.05cm}]
  % raias e seus títulos
  \foreach \i/\t in {0/Sanepar, 1/Município, 2/Comunidade} {
    \node[raia=15.6cm, minimum height=1.9cm] at (0,{-\i*1.95}) {};
    \node[titulo raia] at (0.9,{-\i*1.95-0.95}) {\t};
  }
  \draw[gray!60] (1.8,0) -- (1.8,-5.8);

  % Sanepar
  \node[passo] (s1) at (3.2,-0.95)  {Levantamento e estudo técnico de engenharia};
  \node[passo] (s3) at (5.9,-0.95)  {Projetos de engenharia do SAA};
  \node[passo] (s4) at (8.6,-0.95)  {Fornecimento de materiais hidráulicos};
  \node[passo] (s8) at (14.0,-0.95) {Monitoramento e suporte pós-obra};
  % Município
  \node[passo] (s2) at (3.2,-2.9)   {Termo de Compromisso e Responsabilidade};
  \node[passo] (s5) at (8.6,-2.9)   {Execução das obras, com apoio técnico da Sanepar};
  % Comunidade
  \node[passo] (s6) at (11.3,-4.85) {Formação da associação de moradores};
  \node[passo] (s7) at (14.0,-4.85) {Treinamento e educação ambiental};

  \draw[seta] (s1) -- (s2);
  \draw[seta] (s2.east) -| (s3.south);
  \draw[seta] (s3) -- (s4);
  \draw[seta] (s4) -- (s5);
  \draw[seta] (s5.south) |- (s6.west);
  \draw[seta] (s6) -- (s7);
  \draw[seta] (s7) -- (s8);
\end{tikzpicture}
